\documentclass{article}
\usepackage[T2A]{fontenc}
\usepackage[cp1251]{inputenc}
\usepackage{amsthm}
\usepackage{amsmath}
\usepackage{amssymb}
\usepackage{amsfonts}
\usepackage{mathrsfs}
\usepackage[12pt]{extsizes}
\usepackage{fancyvrb}
\usepackage{indentfirst}
\usepackage[
  left=2cm, right=2cm, top=2cm, bottom=2cm, headsep=0.2cm, footskip=0.6cm, bindingoffset=0cm
]{geometry}
\usepackage[english,russian]{babel}

\begin{document}
\section*{Вариант 19}

Упорядоченный группоид $(A,\cdot,\le)$ тогда и только тогда принадлежит классу $R\{*,\subset\}$, когда для всякой $n$-диады $\omega=(\alpha,\beta)$ и любых термов $p_1,\ldots, p_n; p_0, p_n$ $(i,j=0,\ldots,l_{n})$ таких, что $G_{k-1}^{(\alpha(k),\beta(k))}<G(p_{k})$ $(k=1,\ldots,n)$, $G_{\omega}^{(r,r+1)}<G(\tilde{p}_{t})$ $(t=0,\ldots,n;$ $r=0$ для $t=0$ и $r=th+1$ для $t=1,\ldots,n)$ и $G_{\omega}^{(i,j)}<G(q_{i,j})$ $(i, j = 0,\ldots, l_n)$, выполняются аксиомы:
\begin{gather*}
\left(\bigwedge_{k=0}^{n}x_{3k+1}\ne0\wedge\bigwedge_{k=1}^{n}p_{k}\le x_{3k-1}x_{3k}\right)\rightarrow x_{3t+1}\le\tilde{p}_{t}, \\
\left(\bigwedge_{k=0}^{n}x_{3k+1}\ne0\wedge\bigwedge_{k=1}^{n}p_{k}\le x_{3k-1}x_{3k}\right)\rightarrow q_{i,j}\ne0,
\end{gather*}
где 0 - нулевой элемент упорядоченного группоида $(A,\cdot,\le)$.

\end{document}